\documentclass[11pt,a4paper]{article}
\usepackage[T1]{fontenc}\usepackage[utf8]{inputenc}\usepackage{lmodern}
\usepackage[a4paper,margin=2.2cm]{geometry}
\usepackage{amsmath,amssymb}\usepackage{enumitem}\usepackage{booktabs,longtable,array}
\usepackage[hidelinks]{hyperref}
\setlength{\parindent}{0pt}\setlength{\parskip}{0.4em}
\title{\textbf{OP-CONST Step 3A/3B: anchoring classification of $Z_A$ and
$\mathcal I_{\rm EW}$}\\\large Recon v1 --- proposal-only; no preprint edits}
\author{Per accepted Step 1--2 audit (rev.~1.1) and review directive}
\date{12 June 2026}
\begin{document}\maketitle

\section*{0. Scope and firewall}
This document executes Step 3A (anchoring class of $Z_A$) and Step 3B (anchoring
class of $\mathcal I_{\rm EW}$) at recon level: it classifies what the recorded
preprint structures imply, derives the quantitative protection requirements, and
defines the Step-4 derivation targets. All bounds remain order-of-magnitude stress
anchors pending the citation pass. No D1 block is drafted; no preprint text is
touched. Required protection precisions used throughout (retained band,
representative channel redshifts): $|\zeta_A|\lesssim(1.5$--$1.8)\times10^{-5}$
(quasar $\alpha$, $z\sim1.5$), $|c_{v_2}|=|\partial\mathcal I_{\rm EW}/\partial\ln
u|\lesssim(1.1$--$1.4)\times10^{-6}$ (H$_2$ $\mu$, $z\sim2.5$).

\section{Step 3A --- the anchoring class of $Z_A$}
\textbf{(a) The recorded composition.} Appendix A.3 of the fine-structure route
records the composition class explicitly: ``$Z_B^{(0)}$ inherits from the
compact-phase stiffness $K_\theta^{\rm eff}\sim\phi_0^2$ \dots\ and $Z_W^{(0)}$
inherits from the orientation-sector stiffness $\kappa_n\sim\phi_0^2$''
(L$\sim$52179--52185), with the orientation stiffness recorded elsewhere as the
induced chain $\kappa_n=\mathcal C_n u_0^2$ (OP-C$_n$, L$\sim$4439). The
gauge-kinetic coefficients are therefore \emph{stiffness-induced} quantities of the
condensate, not primitive inputs.

\textbf{(b) Consequence for the dichotomy.} This upgrades the status of the $p=2$
case: it is no longer merely a borrowed analogy but the \emph{face-value reading of
the recorded composition} --- both $Z^{(0)}$ coefficients are written as
$\sim(\text{scale})^2$ stiffnesses of exactly the induced type. The entire Step-3A
question therefore reduces to the evaluation point of these stiffnesses:
P-anchored (the parameter combination $\phi_0=\sqrt{\mu^2/\lambda}$, frozen) versus
B-tracking (the local/cosmological background amplitude $u(X)$, which the retained
$\varepsilon$-band asserts to drift). \emph{Nothing in the current record
establishes the P-anchored reading; the burden of proof is now on protection.}

\textbf{(c) Two distinct $u$-entry points.} The recorded route exposes two
independent places where background drift could enter $\alpha$:
(i)~the microscopic stiffnesses $K_\theta^{\rm eff},\kappa_n$ inside
$(Z_B^{(0)},Z_W^{(0)})$ --- the power-law channel, $p\simeq2$ at face value;
(ii)~the matching-interval endpoint of the recorded one-loop running,
$\ln(\phi_0/M_Z)\approx37.8$: if the upper endpoint tracks $u$ rather than the
parameter-level $\phi_0$, then $\delta\alpha^{-1}=-(b_{\rm eff}/2\pi)\,\delta\ln u$,
a logarithmic channel with $c_\alpha=\alpha\,b_{\rm eff}/2\pi$.

\textbf{(d) Scenario ladder (stress layer; representative $z=1.5$,
$\varepsilon_*$).}

\renewcommand{\arraystretch}{1.15}\small
\begin{longtable}{@{}p{5.4cm}p{2.6cm}p{2.4cm}p{3.0cm}@{}}
\toprule \textbf{Scenario} & $|c_\alpha|$ & $|\Delta\alpha/\alpha|$ &
\textbf{vs $10^{-6}$ anchor} \\ \midrule \endhead \bottomrule \endfoot
$p=2$ stiffness tracking (face-value) & $2$ & $0.12$ & $\sim10^{5}$ tension \\
$p=1$ & $1$ & $0.06$ & $\sim6\times10^{4}$ \\
endpoint-log, $b_{\rm eff}=b_Y=41/6$ & $7.9\times10^{-3}$ & $4.7\times10^{-4}$ &
$\sim470$ \\
endpoint-log, $b_{\rm eff}=1$ & $1.2\times10^{-3}$ & $6.9\times10^{-5}$ &
$\sim70$ \\
required for safety & $\lesssim1.7\times10^{-5}$ & $\lesssim10^{-6}$ & --- \\
\end{longtable}\normalsize
Even the mildest natural tracking channel (endpoint-log) sits two orders above the
anchor; protection must therefore hold at \emph{both} entry points.

\textbf{(e) Protection-mechanism candidates (Step-4 derivation targets, not
claims).}
(P1)~\emph{Normalisation-ratio cancellation}: the physical coupling may depend only
on a ratio in which the stiffness cancels --- e.g.\ if the emergent gauge field is
itself normalised by the same stiffness, $q^2/Z$ could reduce to a parameter-only
combination. Requires the explicit $U(1)$-emergence normalisation
(\S sec:u1\_emergence) to be computed.
(P2)~\emph{IR re-anchoring}: the Thomson-limit $\alpha$ is fixed by IR data; if all
$u$-dependence resides above thresholds that themselves track $u$, the IR value
could be invariant. Requires the threshold bookkeeping of A.4 to be made
$u$-explicit.
(P3)~\emph{Compactness theorem}: the compact-$U(1)$ periodicity that quantises the
charge lattice might also fix the stiffness normalisation to a parameter-level
quantity. Requires the compactness construction to be re-examined for
normalisation, not only quantisation.

\textbf{(f) Free corollary (ratio protection of $\theta_W$).} Since
$\sin^2\theta_W(\phi_0)=Z_W/(Z_B+Z_W)$ (recorded), \emph{any} common tracking
$Z_B,Z_W\propto u^{p}$ with equal $p$ cancels in the ratio: the weak mixing angle
is drift-protected at leading order even in the worst $\alpha$-scenario, up to the
recorded structural difference between the two stiffnesses. This is a clean,
cheap partial protection statement and a candidate formal lemma for Step 4.

\section{Step 3B --- the anchoring class of $\mathcal I_{\rm EW}$}
\textbf{(a) Recorded status.} $\mathcal I_{\rm EW}\equiv\ln(\phi_0/v_2)\approx36.8$
(eq:I\_EW\_definition) is a \emph{definition} from matched values, not a mechanism;
the preprint records four candidate templates, none derived from P1--P6, with the
phase--orientation alignment route as the only structural candidate. The derivative
$\partial\mathcal I_{\rm EW}/\partial\ln u$ is therefore
\emph{mechanism-dependent}, and Step 3B becomes a sensitivity audit of the four
recorded templates.

\textbf{(b) Template-by-template anchoring sensitivity.}
\begin{enumerate}[leftmargin=1.9em,itemsep=2pt,label=(\roman*)]
\item \emph{Asymptotic-freedom-like}: $\mathcal I_{\rm EW}=8\pi^2/[b_a
g_a^2(M_{\rm end})]$. With $dg_a^2/d\ln M=-(b_a/8\pi^2)g_a^4$ one gets exactly
\[
\frac{\partial\mathcal I_{\rm EW}}{\partial\ln M_{\rm end}}=+1 .
\]
The template is \emph{binary}: P-anchored endpoint ($M_{\rm end}=\phi_0$ as a
parameter) gives $0$; a $u$-tracking endpoint gives the full-strength
$|c_{v_2}|=1$, in $\sim4\times10^{5}$ tension with the $\mu$ anchor.
\item \emph{Walking / near-conformal}: $\mathcal I_{\rm EW}\sim\Delta\eta/
\beta_{\rm eff}$ --- a ratio of fixed-point data. Fixed-point quantities are
parameter-level by nature: this template is a \emph{naturally P-anchored
candidate} (subject to the recorded $\gamma_m\gtrsim1$ protection requirement).
\item \emph{Miransky/BKT}: $\mathcal I_{\rm EW}=\pi/\sqrt\delta$ with
$\delta\approx7.3\times10^{-3}$ the distance to the critical surface. Then
\[
\frac{\partial\mathcal I_{\rm EW}}{\partial\ln u}
=-\frac{\mathcal I_{\rm EW}}{2\delta}\,\frac{\partial\delta}{\partial\ln u}
\approx-2.5\times10^{3}\,\frac{\partial\delta}{\partial\ln u},
\]
an \emph{amplified} sensitivity: safety requires
$|\partial\delta/\partial\ln u|\lesssim5\times10^{-10}$. If the critical distance
has any background component, this template is hyper-exposed.
\item \emph{Topological / Hopf-like}: $\mathcal I_{\rm EW}=S_{\rm top}$, a
dimensionless (quantised or geometric) action: \emph{P-anchored by construction}
if genuinely topological.
\end{enumerate}

\textbf{(c) Selection principle (new export to OP-EW-scale).} The varying-constants
channels impose $|\partial\mathcal I_{\rm EW}/\partial\ln u|\lesssim1.2\times
10^{-6}$ on the retained band. The four templates carry anchoring sensitivities
$\{1;\ \approx0\ (\text{fixed-point data});\ 2.5\times10^{3}\cdot\partial_{\ln
u}\delta;\ 0\ (\text{topological})\}$. \emph{The filter favours the fixed-point
(ii) and topological (iv) classes, makes (i) binary in its endpoint anchoring, and
renders (iii) viable only with hyper-precise anchoring of the critical distance.}
Any future OP-EW mechanism must pass this filter --- a new structural selection
criterion exported from OP-CONST to the OP-EW programme.

\section{Combined verdict and stakes}
\textbf{Outcome A} (structural protection) now requires actual theorems, not
default readings: for $\alpha$, a protection mechanism of class (P1)--(P3) against
the face-value stiffness composition; for $\mu$, restriction of the
$\mathcal I_{\rm EW}$ mechanism to the P-anchored classes (ii)/(iv) or a proved
P-anchored endpoint in class (i). \textbf{Outcome B} (screening) remains
unavailable for the cosmological channels, as established at Step 1--2.
\textbf{Outcome C} (limitation) is quantified: any $O(1)$ tracking of either
quantity forces $\varepsilon\lesssim5.5\times10^{-7}$, i.e.\ the retained
five-probe band $[0.0296,0.0376]$ would be in tension by a factor
$\sim5\times10^{4}$ --- the $\varepsilon$-sector would not survive in its present
form. Honest status line: \emph{nothing in the current record establishes the
protection; this is now the sharpest open consistency requirement of the theory,
and it was found internally rather than by referees.}

\section{Step 4 definition (derivation targets; proposal-layer, review-gated)}
\begin{enumerate}[leftmargin=1.9em,itemsep=2pt,label=\textbf{T\arabic*.}]
\item Compute the $U(1)$-emergence normalisation explicitly: does the physical
coupling reduce to a parameter-only combination $q^2/Z\to f(\mu^2,\lambda,
\alpha,\beta)$ (P1), or does the stiffness survive in the observable (B-tracking
confirmed)? This is open target (vii) of the preprint, now carrying a quantitative
acceptance threshold $|\zeta_A|\lesssim1.7\times10^{-5}$.
\item Check the $\sigma$-mode/amplitude cancellation: induced stiffnesses and
induced masses both scale with $u$; determine whether their recorded ratios cancel
the tracking in $Z^{(0)}$.
\item Make the threshold bookkeeping of Appendix A.4 $u$-explicit (P2), including
the $\ln(\phi_0/M_Z)$ endpoint classification.
\item Formalise the $\theta_W$ ratio-protection lemma (Step 3A(f)).
\item Formalise the charge-ratio protection theorem (Step 1--2, F5).
\item For OP-EW: record the selection principle of \S2(c) as a constraint on all
four templates; for template (iii), derive or bound
$\partial\delta/\partial\ln u$.
\end{enumerate}
D1 candidates after review of T1--T3: an OP-CONST problem-block entry (locking
matrix + protection requirements + the selection principle) and a one-sentence
constraint in the OP-EW discussion. Both remain deferred under the recon firewall.
\end{document}
